%! Multiplying \& Dividing Radicals; Rationalizing — Unit 6, Lesson 6.3 — Exit Ticket (Answer Key)
\documentclass[10pt]{article}
\usepackage{algebra2-article}
\usepackage{algebra2-key}   % pulls in -boxes; do NOT also load -boxes
\newcommand{\TallMath}[1]{$\displaystyle #1\rule[-1.4em]{0pt}{3.2em}$}

\begin{document}

\pageheader{Unit 6, Lesson 6.3}{Exit Ticket --- Answer Key}
\namedateperiod

\vspace{0.10in}

No notes, no calculator. Every answer must be in simplest radical form, with no radical left in a
denominator.

\begin{enumerate}[leftmargin=1.8em, itemsep=10pt]
  \item \textbf{Simplify completely.}

        \vspace{4pt}
        (a) $\sqrt{6}\cdot\sqrt{18} = $ \ans{$6\sqrt{3}$} \hspace{0.4in}
        (b) $\left(2+\sqrt{5}\right)\left(4-\sqrt{5}\right) = $ \ans{$3+2\sqrt{5}$}

        \vspace{6pt}
        (c) $\dfrac{6}{\sqrt{3}} = $ \ans{$2\sqrt{3}$} \hspace{0.55in}
        (d) $\dfrac{2}{\sqrt{7}-\sqrt{5}} = $ \ans{$\sqrt{7}+\sqrt{5}$}

        \vspace{2pt}
        {\small\emph{Work:} (a) $\sqrt{108}=\sqrt{36\cdot3}$ --- neither factor simplified alone,
        the product did. \; (b) $8-2\sqrt5+4\sqrt5-5$. \; (c) $\frac{6\sqrt3}{3}$. \;
        (d) conjugate $\sqrt7+\sqrt5$; denominator $7-5=2$, which cancels the $2$ on top.}

  \item \textbf{Explain.} A classmate says that rewriting $\dfrac{1}{\sqrt{2}}$ as
        $\dfrac{\sqrt{2}}{2}$ ``changes the number.'' Show that the two expressions are equal, and
        explain why multiplying by $\dfrac{\sqrt{2}}{\sqrt{2}}$ is always safe.
        \ansline{Both are $\approx0.7071$ ($1\div1.414$ and $1.414\div2$). Algebraically,
        $\frac{1}{\sqrt2}\cdot\frac{\sqrt2}{\sqrt2}=\frac{\sqrt2}{2}$, since $\sqrt2\cdot\sqrt2=2$.
        It is safe because $\frac{\sqrt2}{\sqrt2}=1$, and multiplying by $1$ leaves the value alone
        --- exactly the way $\frac34=\frac{15}{20}$. Only the \emph{form} changed: the radical moved
        from the bottom to the top. Rationalizing is a rewrite, not a new number.}

  \item \textbf{SOL-style multiple choice.} Which expression is equivalent to
        $\dfrac{4}{2-\sqrt{3}}$?
        \begin{enumerate}[label=\Alph*., leftmargin=1.9em, itemsep=1pt]
          \item $\dfrac{8+4\sqrt{3}}{7}$
          \item $8-4\sqrt{3}$
          \item $8+4\sqrt{3}$ \quad \textcolor{keyred}{\textbf{$\leftarrow$ correct}}
          \item $2+\sqrt{3}$
        \end{enumerate}
        \vspace{2pt}
        \begin{itemize}[leftmargin=1.4em, nosep]
          \item \textbf{C} is correct: multiply by $\dfrac{2+\sqrt3}{2+\sqrt3}$. The denominator is
                $4-3=1$, so the answer is $4\left(2+\sqrt3\right)=8+4\sqrt3$.
          \item \textbf{A} used the conjugate but \emph{added} in the denominator ($4+3=7$) instead
                of subtracting --- the difference-of-squares sign.
          \item \textbf{B} multiplied by $2-\sqrt3$ (the same binomial) rather than the conjugate,
                then reported a numerator that never had a rational denominator under it.
          \item \textbf{D} divided the $4$ away, as though $\dfrac{4}{2-\sqrt3}$ reduced by
                cancelling the $4$ with the $2$.
        \end{itemize}
\end{enumerate}

\begin{teachernote}
\small Graded for completion. Four piles, which decide how much of 6.4 opens with algebra.
\begin{itemize}[leftmargin=1.3em, nosep]
  \item \textbf{Got it} --- 1(a)--(d) right and item 3 = C.
  \item \textbf{Did not simplify} --- 1(a) left as $\sqrt{108}$. One prompt: ``is that
        \emph{simplest}?''
  \item \textbf{Middle-term droppers} --- 1(b) answered $3$. Point at Warm-Up 1(c); they rejected
        $(x+4)^2=x^2+16$ an hour ago.
  \item \textbf{Conjugate confusion} --- item 3 = A or B, or 1(d) multiplied by $\sqrt7-\sqrt5$
        again. The pile that matters: A is a sign slip, B is the real misunderstanding. Re-run
        $\left(4+\sqrt3\right)\left(4-\sqrt3\right)=13$ beside $\left(4+\sqrt3\right)^{2}
        =19+8\sqrt3$ --- the contrast is the argument.
\end{itemize}
\textbf{Item 2 is the conceptual item.} A student who cannot say \emph{why} rationalizing is legal
treats it as a ritual and will not reach for it in 6.5, where nothing prompts for it by name.
Accept anything naming $\frac{\sqrt2}{\sqrt2}=1$.
\end{teachernote}

\end{document}
